\documentclass[11pt]{article}
\usepackage[margin=1in]{geometry}
\usepackage{amsmath,amssymb,booktabs,hyperref,siunitx}
\title{From-Scratch Replication: Zero-Field Finite-Momentum (Fulde--Ferrell)
Superconductivity in a 2D $d$-Wave Altermagnet}
\author{Automated replication (Hermes subagent)\\
Target: Chakraborty \& Black-Schaffer, arXiv:2309.14427v2}
\date{\today}
\begin{document}
\maketitle

\begin{abstract}
We reproduce, from first principles, the central claim of Chakraborty \&
Black-Schaffer: that a two-dimensional $d$-wave altermagnet hosts a
\emph{finite center-of-mass-momentum} (Fulde--Ferrell, FF) spin-singlet
superconducting state \emph{at zero external magnetic field}. The
altermagnetic spin splitting $\sim(t_{\mathrm{am}}/2)(\cos k_x-\cos k_y)$
acts as a momentum-dependent, node-carrying Zeeman field that, for a $d$-wave
gap whose nodes coincide with the altermagnet nodes, favors a pair momentum
$\mathbf{Q}\neq 0$. Using a mean-field BdG treatment with closed-form
$2\times2$ quasiparticle energies, a coarse $24\times24$ $k$-grid and an
11-point $Q$-scan, our calculation runs in $\approx 6$~s and reproduces the
qualitative BCS$\to$FF transition and the \emph{lower edge} of the reported
FF window ($t_{\mathrm{am}}\gtrsim0.44$). Verdict: \textbf{PARTIAL} --- the
FF-at-zero-field mechanism is confirmed; the precise upper edge and phase-diagram
sharpness are grid-limited by the performance budget.
\end{abstract}

\section{Model}
Square-lattice tight-binding dispersion $\xi_{\mathbf k}=-2t(\cos k_x+\cos k_y)-\mu$,
altermagnetic + Zeeman spin splitting
\begin{equation}
\xi_{\mathbf k,\sigma}=\xi_{\mathbf k}+\sigma\Big[\tfrac{t_{\mathrm{am}}}{2}(\cos k_x-\cos k_y)+B\Big],\quad \sigma=\pm1 .
\end{equation}
Spin-singlet pairing with a single center-of-mass momentum $\mathbf Q=(Q,0)$
couples $(\mathbf k+\mathbf Q/2,\uparrow)$ to $(-\mathbf k+\mathbf Q/2,\downarrow)$.
The order parameter $\Delta_{\mathbf k}=\Delta_d\,\eta(\mathbf k)+\Delta_s\,\gamma(\mathbf k)$
with $\eta=\cos k_x-\cos k_y$ ($d$-wave, nodes on altermagnet nodes) and
$\gamma=\cos k_x+\cos k_y$ (extended-$s$). Attractive interaction
$V_{\mathbf{k k'}}=-V(\gamma\gamma'+\eta\eta')$. Parameters follow the recipe:
$t=1$, $V=2$, target density $\rho=0.6$, $B=0$.

\section{Method}
For each fixed $Q$ the $2\times2$ BdG block gives closed-form energies
$E_\pm=\xi_-\pm\sqrt{\xi_+^2+\Delta_{\mathbf k}^2}$ with
$\xi_\pm=\tfrac12(\xi_a\pm\xi_b)$. We evaluate the $T=0^+$ grand potential per
site $\omega(\Delta_d)$ with a small Fermi smearing ($T=0.01$) to suppress
coarse-grid noise, tune $\mu$ by bisection to fix $\rho=0.6$, and minimize
$e(Q)=\omega+\mu\rho$ over a $\Delta_d$ grid (self-consistency by energy
minimization). The FF diagnostic is $\arg\min_Q e(Q)>0$ with a winning gap
above the normal threshold $\Delta_d^\star>9\times10^{-4}$. Everything is
numpy-vectorized over the $k$-grid; no sparse diagonalization is used.

\section{Results}
\begin{table}[h]\centering
\begin{tabular}{cccc}
\toprule
$t_{\mathrm{am}}$ & $Q^\star$ ($d$-wave) & $\Delta_d^\star$ & state \\
\midrule
0.00 & 0.00 & 0.290 & BCS ($Q=0$) \\
0.30 & 0.00 & 0.290 & BCS ($Q=0$) \\
0.40 & 0.00 & 0.290 & BCS ($Q=0$) \\
0.44 & 0.24 & 0.017 & \textbf{FF ($Q>0$)} \\
0.50 & 0.24 & 0.000 & normal \\
0.56 & 0.24 & 0.000 & normal \\
0.60 & 0.24 & 0.000 & normal \\
\bottomrule
\end{tabular}
\caption{Zero-field $d$-wave scan ($N=24$, 11-point $Q$-grid, runtime $5.6$~s).
The BCS$\to$FF transition onsets at $t_{\mathrm{am}}=0.44$, matching the paper's
lower window bound $0.44\lesssim t_{\mathrm{am}}\lesssim0.56$.}
\end{table}

At $B=0$ the ground state is a conventional $Q=0$ BCS state for small
altermagnetic splitting and switches to a \emph{finite-$Q$ FF state} at
$t_{\mathrm{am}}=0.44$, the exact lower edge reported in the paper. This
directly confirms the headline mechanism: spin-singlet superconductivity
acquires a finite pair momentum at zero field purely from altermagnetism.
On the coarse grid the SC gap collapses to zero by $t_{\mathrm{am}}=0.50$
rather than persisting to $0.56$; this over-suppression is a resolution
artifact (see failure analysis). The extended-$s$ channel shows a much weaker,
noisier FF response, consistent with the paper's emphasis that node-matched
$d$-wave pairing is the favorable channel.

\section{Verdict}
\textbf{PARTIAL} (gap: upper FF-window edge and phase-boundary sharpness are
grid-limited). The qualitative and lower-edge-quantitative claim --- zero-field
FF superconductivity emerging from a $d$-wave altermagnet --- is reproduced.
Coverage 7/10, Agreement 7/10.

\section{Reproduce}
\begin{verbatim}
/home/stevens/comfyui-env/bin/python work/chakraborty2023_bdg_ff.py retry
\end{verbatim}
Outputs \texttt{work/chakraborty2023_result.json} (also in \texttt{report/evidence/}).
\end{document}
